% Source: MQ34a_v1.html
\documentclass[11pt]{article}
\usepackage[utf8]{inputenc}
\usepackage[T1]{fontenc}
\usepackage[margin=1in]{geometry}
\usepackage{amsmath}
\usepackage{amssymb}
\usepackage{siunitx}
\usepackage{enumitem}
\usepackage{parskip}

\title{PS160 -- MQ34a\_v1.html}
\date{}

\begin{document}
\maketitle

\section*{Questions}

\begin{enumerate}[leftmargin=*]
  \item An object is placed 21 cm to the left of a flat mirror.

\textbf{(a)} What is the position of the image? (Positive on the same side as the object; negative on the opposite side. Answer with an integer.) [BLANK]
cm

\textbf{(b)} What is the magnification of the image? (Answer with an integer.) [BLANK]

\textbf{(c)} State whether the image is real or virtual. [BLANK]

\textbf{(d)} State whether the image is upright or inverted. [BLANK] \hfill \textit{(1 pts, Fill-in-blank)}

  \item A spherical mirror is concave and has focal length 11.8 cm. If the object distance is 8.4 cm, calculate the image distance, \emph{s}', in cm. \hfill \textit{(1 pts, Numeric)}

  \item A convex mirror has a focal length of 12 cm. If the object is placed 20 cm from the mirror, the image is found at [DROPDOWN: 7.5 | -5.5 | -7.5 | 15 | -18]
cm, the magnification is [DROPDOWN: -5.25 | 5.25 | -0.375 | 0.375 | -2.67 | 2.67]
, the image is [DROPDOWN: real | virtual]
and [DROPDOWN: upright | inverted]
. \hfill \textit{(1 pts, Dropdown)}

  \item A thin converging lens has focal length 10 cm. If the object distance is 19.1 cm, calculate the image distance in cm. \hfill \textit{(1 pts, Numeric)}

  \item A thin diverging lens has a focal length of 9.00 cm. If the object distance is 6.0 cm, the image distance is [DROPDOWN: -3.6 | -5.8 | 5.8 | 3.6 | -6.6 | 6.6]
cm, the magnification is [DROPDOWN: -1.67 | 0.60 | -0.6 | 3.0 | 1.67]
, and the image is both [DROPDOWN: real | virtual]
and [DROPDOWN: upright | inverted]
. \hfill \textit{(1 pts, Dropdown)}

\end{enumerate}

\end{document}